\documentclass[10pt,a4paper]{amsart}
\usepackage[margin=19mm]{geometry}
\usepackage{amsmath,amssymb,mathtools}
\usepackage[scr=rsfs,cal=euler]{mathalfa}
\usepackage{xfrac}
\usepackage[unicode,hidelinks,bookmarks=false]{hyperref}
\newcommand{\ie}{i.e.\ }
\newcommand{\cf}{cf.\ }
\newcommand{\resp}{respectively\ }
\newcommand{\Subsection}{\subsection}
\providecommand{\nobreakdash}{\nobreak\hbox{-}}
\setlength{\emergencystretch}{2em}
\multlinegap0pt
\usepackage{amssymb}
\usepackage{enumerate}
\newtheorem{theorem}{Theorem}
\newcommand{\orG}{\ensuremath{\overrightarrow{G}}}
\title{Arc-distinguishing of orientations of graphs}
\author{Aleksandra Gorzkowska, Jakub Kwa\'sny}
\date{Source: arXiv:2402.16169v2 (CC BY 4.0)}
\begin{document}
\maketitle

\noindent\emph{Source and licence.} Arc-distinguishing of orientations of graphs. Version arXiv:2402.16169v2 (CC BY 4.0), distributed by arXiv under the Creative Commons Attribution 4.0 International licence, http://creativecommons.org/licenses/by/4.0/, as recorded in the arXiv licence metadata for this exact version and shown on the abstract page https://arxiv.org/abs/2402.16169v2. Full licence notice: \texttt{licenses/arxiv-2402.16169-CC-BY-4.0.txt}.\par\smallskip

\noindent\textit{Editorial scope.} Counted unit: the article's theorem bipartite\_bounds (source0.tex lines 159--170). The article's complete original proof of it is delivered with it (source0.tex lines 171--190, one \texttt{proof} environment of 20 lines). No mathematics is written, repaired or re-derived: the statement and the proof are the article's own lines, in the article's own notation, and no retained line is substituted or re-flowed.  Retained uncounted as the source-local context the proof needs: lines 106--108.  Nothing was omitted from any retained span, so the package carries no omission record.  No retained line contains a citation, so the wrapper carries no bibliography and no bibliography entry.  No retained line contains a figure environment, an image-inclusion command or drawing code, so no image is delivered and the package declares no asset dependency.  Editorial formatting only, in addition: an AMS \texttt{amsart} preamble that restates the article's own declarations and macros used by the retained lines, namely the environments \texttt{theorem} on separate counters, together with the standard packages those lines need.  The retained environments print in this document as Theorem 1 in order of appearance, whereas the article prints the same items with the numbering of its own full document.  The complete native source of the article as distributed in the arXiv e-print for this exact version is bundled unchanged as \texttt{sources/155106/source0.tex}.
\begin{theorem} \label{lemma: part_indep}
Let $G=(V,E)$ be a graph with a partition $V=V_1\cup\dots\cup V_k$ into $k\geq 1$ independent sets, each setwise fixed by every automorphism of $G$. Then $OD'^+(G)=D'(G)$ and $OD'^-(G)=\lceil D'(G)/2\rceil$.
\end{theorem}

\begin{theorem} \label{thm: bipartite_bounds}
Let $G$ be a connected bipartite graph of order at least three. Then
\[
\left\lfloor\frac{D'(G)}{2}\right\rfloor
\leq OD'^-(G)
\leq \left\lfloor\frac{D'(G)}{2}\right\rfloor+1
\]
and
\[
D'(G)-1\leq OD'^+(G)\leq D'(G).
\]
\end{theorem}

\begin{proof}
Let $G=(X\cup Y,E)$. Since $G$ is connected, every automorphism either preserves both bipartition classes setwise or interchanges them.

We first note that any edge colouring with $t$ colours that breaks all non-trivial automorphisms preserving both classes implies $D'(G)\leq t+1$. Indeed, at most one non-trivial automorphism can preserve such a colouring. Otherwise, two distinct colour-preserving automorphisms $\varphi$ and $\psi$ would both interchange the classes, so $\varphi^{-1}\psi$ would be a non-trivial colour-preserving automorphism preserving both classes, a contradiction. If a non-trivial colour-preserving automorphism remains, it moves some edge $e$, since $G$ is connected and has at least three vertices. Recolour $e$ with a new colour. Every automorphism preserving the new colouring fixes $e$ setwise and also preserves the original colouring, so it must be the identity. If the original colouring is already distinguishing, no additional colour is needed.

To prove the lower bound for $OD'^-(G)$, take any orientation $\orG$ with a distinguishing arc colouring using $k$ colours. Apply the same encoding of directions and arc colours as in the proof of Theorem \ref{lemma: part_indep}, taking $V_1=X$ and $V_2=Y$. This gives an edge colouring with at most $2k$ colours that breaks all non-trivial automorphisms preserving both bipartition classes: any such automorphism preserving the encoded colouring would also preserve the orientation and its distinguishing arc colouring. The preceding observation gives $D'(G)\leq 2k+1$. Taking $k=OD'^-(G)$ yields the desired lower bound.

For the upper bound, put $d=D'(G)$ and take a distinguishing edge colouring with colours $1,\ldots,d$. Keep colour $d$ in a singleton group, and partition the remaining colours into pairs, with one additional singleton if necessary. Replace the colours in each group by a common arc colour. For each pair, orient the edges of its first colour from $X$ to $Y$ and those of its second colour from $Y$ to $X$. Orient the edges in every singleton group from $X$ to $Y$. This uses
\[
1+\left\lceil\frac{d-1}{2}\right\rceil
=\left\lfloor\frac{d}{2}\right\rfloor+1
\]
arc colours. An automorphism preserving this arc colouring cannot interchange $X$ and $Y$, since the non-empty class corresponding to the singleton group $\{d\}$ contains only arcs from $X$ to $Y$. Once the bipartition classes are preserved, the original edge colour can be recovered from the arc colour and its direction. Thus such an automorphism preserves the original distinguishing colouring and is the identity.

Finally, orient every edge from $X$ to $Y$, obtaining an orientation $\orG_0$ whose automorphisms are precisely the automorphisms of $G$ preserving both bipartition classes. A distinguishing arc colouring of $\orG_0$ therefore breaks all non-trivial automorphisms of $G$ preserving both classes. Applying the initial observation with $t=D'(\orG_0)$ gives
\[
D'(G)-1\leq D'(\orG_0)\leq OD'^+(G).
\]
The reverse bound $OD'^+(G)\leq D'(G)$ follows because every distinguishing edge colouring of $G$ also distinguishes each of its orientations.
\end{proof}

\end{document}
